% page 429 du fac-simile (image p-428 du scan)
% bloc 172.7 x 242.0 mm ; marge gauche 29.3 mm
\pgc{7.620mm}{0.000mm}{
\l{\cel{54}421}
\l{}
\l{\sou{paraboloido}.\cel{2}(\sou{geom.)}\cel{4}Qua havas la formo di parabolo. DEFIS.}
\l{}
\l{\sou{paradar}.\cel{2}(\sou{netrans}.)(\sou{aludante armeo-trupi}).\cel{2}Defilar lor}
\l{\cel{3}revueso. - DEF .}
\l{}
\l{\sou{paradigmo}. (\sou{gram.})\cel{2}Tipo di deklino, di konjugo.\cel{2}- DEFIRS.}
\l{}
\l{\sou{paradizo}. I. (\sou{kristanismo)}\cel{2}Loko delicoza, gardeno ube Deo poza-}
\l{\cel{3}bis la viro dat-unesma. - II. (\sou{metaf}.)Remdsyo di la beateso}
\l{\cel{3}cielala. -- Loko, regiono, stando, extreme agreabla. - DEFIRS.}
\l{}
\l{\sou{paradoxo}. Opiniono kontrea a la opiniono komuna. - DEFIRS.}
\l{}
\l{\sou{parafo}. - I. Plumo-stroko diversa-forma quan onu adjuntis a\cel{1}\sou{sua}}
\l{\cel{3}nomo por distingar sua signaturo. - II. Signaturo kurtigura}
\l{\cel{3}quan onu trasas en la marjino di akti por aprobar la sur-}
\l{\cel{3}strekizuri. - DEFI.}
\l{}
\l{\sou{parafino.}\cel{2}\sou{(meio)}. Reziduo de la distilo di la oleo minerala,}
\l{\cel{3}substanco grasa, blanka, kristalatra, uzebla por lumifo. -}
\l{\cel{3}DEFIRS.}
\l{}
\l{\sou{parafrazar}. (\sou{trans.})\cel{2}Developar texto por explikar lu. - DEFIRS.}
\l{}
\l{\sou{paragrafo}.Segmento di chapitro. - DEFIRS.}
\l{}
\l{\sou{paralaxo}\cel{2}(\sou{astron.})\cel{2}Angulo formacita, ye la centro di astro,}
\l{\cel{3}da du rekti qui abutas a du observanti lokizita en du punti}
\l{\cel{3}interdiferanta. - DEFIRS .}
\l{}
\l{\sou{paralela}. (\sou{geom.})\cel{2}(\sou{aludante rekti, surfaci}). Di qui la inter-}
\l{\cel{3}disto esas egala irge-punte.\cel{2}- DEFIRS.}
\l{}
\l{\sou{paralelepipedo}. (\sou{geom.})\cel{2}Solido qua finas per sis paralelogrami}
\l{\cel{3}inter-egala, qui inter-opozas duope. - DEFIS.}
\l{}
\l{\sou{paralelogramo}.- (\sou{geom}.) Quadrilatero di qua la lateri inter-opozata}
\l{\cel{3}esas egala e paralela. - DEFIRS.}
\l{}
\l{\sou{paralizar}. (\sou{trans.}) I. (\sou{patol.}) Produktar la nepovo, tota o parta,}
\l{\cel{3}pri sive movar e sentar, sive movar o sentar. - II. (\sou{metaf.})}
\l{\cel{3}Impedar la funciono normala. - DEFIRS.}
\l{}
\l{\sou{paralogismo}. Rezono nekorekta. - DEFIS.}
\l{}
\l{\sou{paramento}. Speco de stofo, simpla od ornita, adjuntita a la extre-}
\l{\cel{3}majo di la maniko di vesto, uniformo, robo. - FIS.}
\l{}
\l{\sou{parametro}. (\sou{matem.}) Konstanto ek la equaciono, ek la konstrukto di}
\l{\cel{3}kurvo, ek surfaco, qua posibligas la determino di lua dimensioni.}
\l{\cel{3}DEFIS.}
\l{}
\l{\sou{parani}mfo. (en\cel{1}\sou{la epoki antiqua}).\cel{2}Amiko di la quik-mariajonto, qua}
\l{\cel{3}helpis lu pri la ceremoniaro di la mariajo. - DEFIS.}
}
